% Simple methods note for the SAFOD active-source DAS tide experiment.
% Compile with: pdflatex solid_earth_tide_model.tex
\documentclass[11pt]{article}

\usepackage[margin=1in]{geometry}
\usepackage[T1]{fontenc}
\usepackage{lmodern}
\usepackage{amsmath,amssymb}
\usepackage{booktabs}
\usepackage{enumitem}
\usepackage[hidelinks]{hyperref}

\newcommand{\dvv}{\delta v/v}
\newcommand{\eps}{\varepsilon}
\newcommand{\Rearth}{R_{\oplus}}
\newcommand{\Mbody}{\mathcal{M}}

\setlength{\parskip}{0.45em}
\setlength{\parindent}{0pt}

\title{\textbf{A Simple Solid--Earth Tide Forward Model for SAFOD}}
\author{Methods note --- SAFOD active-source borehole DAS experiment}
\date{\today}

\begin{document}
\maketitle

\begin{abstract}
This note explains a first-pass model of the solid-Earth tide at the San
Andreas Fault Observatory at Depth (SAFOD).  The model calculates the changing
gravitational forcing from the Moon and Sun, converts that forcing into a
scalar Earth-deformation proxy using degree-2 Love numbers, and uses the
resulting time series as a template for testing an apparent seismic velocity
change.  The model predicts the astronomical forcing; it does not by itself
predict how strongly SAFOD rock, fractures, fluids, or the DAS measurement
should respond.  That response is represented by an unknown strain sensitivity.
\end{abstract}

\tableofcontents

% ---------------------------------------------------------------------------
\section{The physical question}
% ---------------------------------------------------------------------------

The Moon and Sun exert a time-varying gravitational force on the solid Earth.
Because the force is not exactly the same at every point on Earth, it produces
a differential deformation: the solid-Earth tide.  The deformation is small,
but it repeats on a predictable astronomical schedule.

The SAFOD experiment measures an apparent fractional seismic velocity change,
written
\begin{equation}
  y(t) \equiv \left(\frac{\delta v}{v}\right)_{\!\mathrm{app}}(t).
\end{equation}

The simplest physical coupling model is
\begin{equation}
  \left(\frac{\delta v}{v}\right)_{\!\mathrm{app}}(t)
  = S\,\eps_{\mathrm{tide}}(t) + \text{drift}(t) + \text{noise}(t),
  \label{eq:coupling}
\end{equation}
where $\eps_{\mathrm{tide}}$ is the predicted tidal strain and $S$ is an
unknown rock sensitivity.  The coefficient $S$ absorbs the response of the
formation and the sensitivity of the particular seismic observable.  It is
not assumed to be one.

The forward model therefore answers:

\begin{quote}
What tidal deformation should be present at SAFOD at each time?
\end{quote}

The matched-filter analysis then asks whether the measured $y(t)$ contains
that predicted shape.

% ---------------------------------------------------------------------------
\section{What the first-pass model is and is not}
% ---------------------------------------------------------------------------

The model below is a useful scalar forcing model.  It is not yet a complete
downhole tidal-strain calculation.  In particular, it does not calculate the
full depth-dependent strain tensor projected onto the actual DAS fiber.  The
scalar result should therefore be described as an \emph{areal-strain proxy} or
\emph{tidal forcing template}, not as a direct measurement of the strain in
the SAFOD formation.

This distinction matters because DAS measures strain or strain rate along a
fiber direction.  A more complete model would calculate the strain tensor
$\eps_{ij}$ and project it onto the local fiber unit vector $n$:
\begin{equation}
  \eps_{\mathrm{fiber}}(t) = n_i\,\eps_{ij}(t)\,n_j.
  \label{eq:fiber_projection}
\end{equation}

The scalar model is nevertheless a sensible starting point because the
astronomical time dependence is the main object being tested first.

% ---------------------------------------------------------------------------
\section{Step 1: define the site and times}
% ---------------------------------------------------------------------------

Use the SAFOD site coordinates
\begin{align}
  \phi &= 35.9746^\circ, &
  \lambda &= -120.5516^\circ,
\end{align}
where $\phi$ is latitude and $\lambda$ is longitude.

For each UTC epoch time, calculate the Julian Date (JD) and define the time
relative to J2000.0 as
\begin{equation}
  d = \mathrm{JD} - 2451545.0.
  \label{eq:j2000}
\end{equation}

The time convention must be explicit.  The Sun and Moon positions are
time-dependent, so an incorrect timezone or an incorrect UTC conversion
shifts the phase of the predicted template.

% ---------------------------------------------------------------------------
\section{Step 2: calculate the Sun and Moon geometry}
% ---------------------------------------------------------------------------

At every epoch, the model needs the right ascension $\alpha$, declination
$\delta$, and geocentric distance $D$ of each body.  A low-precision ephemeris
can be used for a short, first-pass experiment.  The Sun can be calculated
from its mean longitude and mean anomaly; the Moon can be calculated from a
truncated orbital solution.  For higher-accuracy work, replace these formulas
with a tested ephemeris or a standard solid-Earth-tide package.

The local hour angle of a body is
\begin{equation}
  H = \mathrm{LST} - \alpha,
  \label{eq:hour_angle}
\end{equation}
where LST is local sidereal time.  The angle $\theta$ between the body and the
local upward direction at SAFOD follows from the spherical triangle:
\begin{equation}
  \cos\theta
  = \sin\phi\,\sin\delta
    + \cos\phi\,\cos\delta\,\cos H.
  \label{eq:zenith_angle}
\end{equation}

The angle changes as the Earth rotates.  The Sun and Moon also have different
declinations and distances, so their contributions are not identical.

% ---------------------------------------------------------------------------
\section{Step 3: calculate the degree-2 tidal potential}
% ---------------------------------------------------------------------------

The dominant body-tide forcing is the degree-2 part of the tidal potential.
For one body $b$ (the Moon or Sun), the simplified potential at the Earth’s
surface is
\begin{equation}
  W_b(t) =
  \frac{G\,\Mbody_b}{D_b(t)^3}\,\Rearth^2\,
  P_2\!\left(\cos\theta_b(t)\right),
  \label{eq:potential}
\end{equation}
where
\begin{equation}
  P_2(x) = \frac{3x^2-1}{2}
  \label{eq:legendre}
\end{equation}
is the second Legendre polynomial.

The total forcing is the sum of the lunar and solar terms:
\begin{equation}
  W(t) = W_{\mathrm{Moon}}(t) + W_{\mathrm{Sun}}(t).
  \label{eq:total_potential}
\end{equation}

The factor $D^{-3}$ is important: tidal forcing depends on the spatial
gradient of the gravitational field, not merely on the gravitational force.
The nearby Moon can therefore make a large tidal contribution despite its
smaller mass.

Because $P_2(\cos\theta)$ is quadratic in the geometry, the result is not
generally a single sinusoid.  The combined signal usually contains diurnal
and semidiurnal components with unequal peak heights.

% ---------------------------------------------------------------------------
\section{Step 4: convert potential into a strain proxy}
% ---------------------------------------------------------------------------

The quantity
\begin{equation}
  \frac{W(t)}{g\Rearth}
  \label{eq:dimensionless_forcing}
\end{equation}
is dimensionless.  It represents the scale of the tidal displacement relative
to the Earth’s radius.

Love numbers describe the elastic response of the Earth to the forcing.  In
the scalar model, $h_2$ is the degree-2 radial-displacement Love number and
$l_2$ is the degree-2 horizontal, or Shida, number.  The first-pass areal
strain is
\begin{equation}
  \boxed{
  \eps_{\mathrm{tide}}(t)
  = \left(2h_2-6l_2\right)\frac{W(t)}{g\Rearth}}
  \label{eq:areal_strain}
\end{equation}

The constants used in the current SAFOD model are:

\begin{center}
\begin{tabular}{@{}lll@{}}
\toprule
Quantity & Meaning & Value \\
\midrule
$R_{\oplus}$ & Earth radius & $6.371\times10^6$ m \\
$g$ & surface gravity & $9.81$ m s$^{-2}$ \\
$h_2$ & radial Love number & $0.6032$ \\
$l_2$ & horizontal Love number & $0.0839$ \\
$G$ & gravitational constant & $6.674\times10^{-11}$ m$^3$ kg$^{-1}$ s$^{-2}$ \\
$\mathcal{M}_{\mathrm{Moon}}$ & lunar mass & $7.342\times10^{22}$ kg \\
$\mathcal{M}_{\mathrm{Sun}}$ & solar mass & $1.989\times10^{30}$ kg \\
\bottomrule
\end{tabular}
\end{center}

For these Love numbers,
\begin{equation}
  2h_2-6l_2 = 0.703.
\end{equation}

Evaluated over the current 24-hour experiment window, this model gives
approximately
\begin{align}
  \eps_{\mathrm{tide}}^{\mathrm{pk-pk}} &\approx 9.09\times10^{-8},\\
  \eps_{\mathrm{tide}}^{\mathrm{amp}} &\approx 4.54\times10^{-8}.
\end{align}

These values describe the predicted forcing scale, not the predicted
$\dvv$.  The conversion from tidal strain to seismic velocity change requires
the unknown sensitivity $S$ in Eq.~\eqref{eq:coupling}.

% ---------------------------------------------------------------------------
\section{Step 5: connect the tide to the measured $\dvv$}
% ---------------------------------------------------------------------------

The measured epoch series is obtained from travel-time changes.  For a fixed
wave path and a small velocity perturbation,
\begin{equation}
  \frac{\delta t}{t} = -\frac{\delta v}{v}.
  \label{eq:travel_time}
\end{equation}

The minus sign means that a faster medium produces an earlier arrival.  In
practice, the measured quantity is called \emph{apparent} $\dvv$ because it
can also contain clock drift, source-wavelet variation, coupling changes, and
processing artifacts.

Let $y_k$ be the apparent $\dvv$ at epoch time $t_k$.  The forward-model curve
is normalized to make a unit-amplitude template:
\begin{equation}
  \tau_k =
  \frac{\eps_{\mathrm{tide}}(t_k)-\overline{\eps_{\mathrm{tide}}}}
       {\max_j\left|\eps_{\mathrm{tide}}(t_j)-\overline{\eps_{\mathrm{tide}}}\right|}.
  \label{eq:template}
\end{equation}

The simplest fit is
\begin{equation}
  y_k = A\tau_k + c_0+c_1t_k+c_2t_k^2+r_k,
  \label{eq:fit_model}
\end{equation}
where $A$ is the fitted apparent-velocity response, the $c$ terms represent
slow drift, and $r_k$ is the remaining residual.

The three useful trend choices are:
\begin{itemize}[nosep]
  \item flat: $c_1=c_2=0$;
  \item linear: $c_2=0$;
  \item quadratic: all three trend coefficients are fitted.
\end{itemize}

The null hypothesis is the trend-only model, $A=0$.  The tide is the one
additional, physically predicted shape being tested.

If the template is scaled using the physical tidal amplitude rather than the
unit normalization in Eq.~\eqref{eq:template}, the strain sensitivity is
\begin{equation}
  S = \frac{A}{\eps_{\mathrm{tide}}^{\mathrm{amp}}}.
  \label{eq:sensitivity}
\end{equation}

For example, a detection floor of $8\times10^{-4}$ in apparent $\dvv$ would
correspond to
\begin{equation}
  S_{\mathrm{floor}}
  = \frac{8\times10^{-4}}{4.54\times10^{-8}}
  \approx 1.8\times10^4.
\end{equation}

This means the rock and measurement system would need an effective sensitivity
of roughly $1.8\times10^4$ for the predicted tidal forcing to reach that
detection floor.

% ---------------------------------------------------------------------------
\section{A compact implementation recipe}
% ---------------------------------------------------------------------------

The forward model can be implemented in the following order:

\begin{enumerate}
  \item Convert every UTC epoch time to Julian Date.
  \item Calculate the Sun’s and Moon’s $\alpha$, $\delta$, and $D$.
  \item Calculate local sidereal time and each body’s hour angle.
  \item Calculate $\cos\theta$ using Eq.~\eqref{eq:zenith_angle}.
  \item Evaluate the degree-2 potential in Eq.~\eqref{eq:potential}.
  \item Sum the solar and lunar potentials.
  \item Convert potential to the scalar strain proxy using Eq.~\eqref{eq:areal_strain}.
  \item Interpolate the model to the measured epoch times.
  \item Fit the normalized template plus flat, linear, or quadratic trend.
  \item Estimate uncertainty using a time-respecting bootstrap or another
        method appropriate for serially correlated epochs.
\end{enumerate}

In pseudocode:

\begin{verbatim}
for t in epoch_times:
    jd = utc_to_julian_date(t)
    sun  = solar_position(jd)
    moon = lunar_position(jd)
    W = 0
    for body in (sun, moon):
        cos_theta = site_body_geometry(body, SAFOD_lat, SAFOD_lon, t)
        W += G * body.mass / body.distance**3 \\
             * R_earth**2 * P2(cos_theta)
    tide_strain[t] = (2*h2 - 6*l2) * W / (g * R_earth)

template = normalize_zero_mean(tide_strain)
fit measured_dvv = A*template + trend + residual
\end{verbatim}

% ---------------------------------------------------------------------------
\section{Limitations and next refinements}
% ---------------------------------------------------------------------------

\begin{enumerate}
  \item \textbf{Ephemeris accuracy.} A low-precision Sun/Moon calculation is
        adequate for understanding the template over one day, but a precise
        study should use a tested ephemeris and a standard body-tide package.

  \item \textbf{Scalar versus tensor strain.} Equation~\eqref{eq:areal_strain}
        is a scalar areal-strain proxy.  A downhole DAS measurement is better
        represented by the strain tensor projected onto the actual fiber
        direction, Eq.~\eqref{eq:fiber_projection}.

  \item \textbf{Surface versus downhole response.} The Love-number expression
        is written as a surface degree-2 response.  The deformation field at
        depth and the response of a deviated borehole may differ from this
        scalar surface approximation.

  \item \textbf{Other loadings.} Ocean tides, atmospheric pressure, hydrology,
        and local poroelastic effects can contribute to deformation.  They are
        separate from the direct Sun--Moon body-tide term.

  \item \textbf{Rock physics.} The astronomical model predicts the forcing,
        not the value of $S$.  A non-detection constrains the product of rock
        sensitivity and measurement quality; it does not prove that the rock
        has zero tidal response.
\end{enumerate}

The first scientifically clean claim is therefore:

\begin{quote}
The Sun--Moon body tide provides a known, small, non-sinusoidal forcing
template at SAFOD.  The active-source DAS experiment tests whether the
measured apparent velocity-change series contains a response with that shape
above its drift and noise floor.
\end{quote}

% ---------------------------------------------------------------------------
\section{References}
% ---------------------------------------------------------------------------

\begin{thebibliography}{99}

\bibitem{Agnew2007}
D.~C. Agnew (2007), ``Earth tides,'' in T.~A. Herring (ed.),
\emph{Treatise on Geophysics}, vol.~3, pp.~163--195, Elsevier.
\href{https://doi.org/10.1016/B978-044452748-6.00056-0}{doi:10.1016/B978-044452748-6.00056-0}.

\bibitem{MetivierConrad2008}
L.~M\'etivier and C.~P. Conrad (2008), ``Body tides of a convecting,
laterally heterogeneous, and aspherical Earth,'' \emph{Journal of Geophysical
Research: Solid Earth}, 113, B11405.
\href{https://doi.org/10.1029/2007JB005448}{doi:10.1029/2007JB005448}.

\bibitem{Farrell1972}
W.~E. Farrell (1972), ``Deformation of the Earth by surface loads,''
\emph{Reviews of Geophysics}, 10(3), 761--797.
\href{https://doi.org/10.1029/RG010i003p00761}{doi:10.1029/RG010i003p00761}.

\bibitem{Poupinet1984}
G.~Poupinet, W.~L. Ellsworth, and J. Fr\'echet (1984), ``Monitoring velocity
variations in the crust using earthquake doublets: An application to the
Calaveras Fault, California,'' \emph{Journal of Geophysical Research}, 89(B7),
5719--5731.
\href{https://doi.org/10.1029/JB089iB07p05719}{doi:10.1029/JB089iB07p05719}.

\bibitem{Takano2014}
T.~Takano, T.~Nishimura, H.~Nakahara, Y.~Ohta, and S.~Tanaka (2014),
``Seismic velocity changes caused by the Earth tide: Ambient noise correlation
analyses of small-array data,'' \emph{Geophysical Research Letters}, 41(17),
6131--6136.
\href{https://doi.org/10.1002/2014GL060690}{doi:10.1002/2014GL060690}.

\bibitem{Lellouch2019}
A.~Lellouch, S.~Yuan, Z.~Spica, B.~Biondi, and W.~L. Ellsworth (2019),
``Seismic velocity estimation using passive downhole distributed acoustic
sensing records: Examples from the San Andreas Fault Observatory at Depth,''
\emph{Journal of Geophysical Research: Solid Earth}, 124, 6931--6948.
\href{https://doi.org/10.1029/2019JB017533}{doi:10.1029/2019JB017533}.

\bibitem{RojstaczerAgnew1989}
S.~Rojstaczer and D.~C. Agnew (1989), ``The influence of formation material
properties on the response of water levels in wells to Earth tides and
atmospheric loading,'' \emph{Journal of Geophysical Research}, 94(B9),
12403--12411.
\href{https://doi.org/10.1029/JB094iB09p12403}{doi:10.1029/JB094iB09p12403}.

\end{thebibliography}

\end{document}
